% Source: https://artofproblemsolving.com/community/c3h2393019_a_question_on_proving
% Stable identity: aops-topic-2393019-q01
\textbf{First part.} $6k+5=6k+3+2=3(2k+1)+2$, so $j=2k+1$ works.

\textbf{The converse fails.} The integer $2=3\cdot0+2$ is of the form $3j+2$. If $2=6k+5$, then $6k=-3$ and $k=-\dfrac12$, which is not an integer. So $2$ is not of the form $6k+5$. (Likewise $8=3\cdot2+2$ is not: $6k+5=8$ gives $k=\dfrac12$.)

\textbf{Exactly which ones.} Suppose $3j+2=6k+5$ for integers $j,k$. Then $3j=6k+3$, so $j=2k+1$ is odd. Conversely, if $j=2k+1$ is odd, then $3j+2=6k+5$. So an integer $3j+2$ is of the form $6k+5$ $\boxed{\text{exactly when } j \text{ is odd}}$; for even $j$ it is not (for example $2$, $8$, $14$).
